\documentclass[11pt]{article}

\usepackage{amsmath,amsthm,amssymb}
\usepackage[margin=1.2in]{geometry}
\usepackage{hyperref}

\newtheorem{theorem}{Theorem}
\newtheorem{corollary}[theorem]{Corollary}
\newtheorem{lemma}[theorem]{Lemma}
\theoremstyle{remark}
\newtheorem{remark}[theorem]{Remark}

\title{Eigenvalue Positivity for the Hecke Staircase Element\\via Kazhdan--Lusztig Positivity}
\author{Clio}
\date{April 16, 2026}

\begin{document}
\maketitle

\begin{abstract}
We prove that all nonzero eigenvalues of the Hecke staircase element
$\Pi_q = \prod_{j=1}^m (1 + T_{i_j})$ acting on any Specht module $V_\lambda$
are positive for all $q > 0$. The proof combines the Elias--Williamson positivity
theorem for Kazhdan--Lusztig basis structure constants with an exterior power
trace argument. As a corollary, the rank of $\Pi_q$ on $V_\lambda$ is constant
for $q > 0$.
\end{abstract}

\section{Setup}

Let $H_q(S_n)$ be the Iwahori--Hecke algebra over $\mathbb{Z}[v, v^{-1}]$
(where $v^2 = q$) with generators $T_s$ satisfying
\[
  T_s^2 = (q-1)\,T_s + q.
\]
Let $C'_s = v^{-1}(1 + T_s)$ be the simple Kazhdan--Lusztig generator.
The \emph{staircase reduced word} for $w_0 \in S_n$ is
$s_{i_1}, \ldots, s_{i_m}$ where $m = \binom{n}{2}$, given explicitly by
\[
  s_1,\; s_2\, s_1,\; s_3\, s_2\, s_1,\; \ldots,\; s_{n-1} \cdots s_1.
\]
Define the \emph{Hecke staircase element}:
\[
  \Pi_q \;=\; \prod_{j=1}^{m} (1 + T_{i_j})
       \;=\; v^m \prod_{j=1}^{m} C'_{s_{i_j}}.
\]

\section{Trace non-negativity}

\begin{theorem}[Trace Non-negativity]\label{thm:trace}
For every partition $\mu \vdash n$, the trace $\operatorname{tr}_\mu(\Pi_q)$
of\/ $\Pi_q$ on the Specht module $V_\mu$ is a polynomial in $q$ with
non-negative integer coefficients.
\end{theorem}

\begin{proof}
\emph{Step~1.}
By the Elias--Williamson positivity theorem~\cite{EW14}, the product of
Kazhdan--Lusztig basis elements has non-negative expansion:
\[
  \prod_{j=1}^m C'_{s_{i_j}}
  \;=\; \sum_{w \in S_n} a_w(v)\, C'_w
  \qquad\text{with } a_w \in \mathbb{Z}_{\geq 0}[v].
\]

\emph{Step~2.}
Let $V_\mu$ be the cell module (Specht module) corresponding to~$\mu$.
In the KL cellular basis $\{\bar{C}_x : x \in \mathrm{cell}(\mu)\}$,
left multiplication by $C'_w$ gives:
\[
  C'_w \cdot \bar{C}_x
  \;=\; \sum_{y \in \mathrm{cell}(\mu)} h_{w,x,y}(v)\, \bar{C}_y
\]
where $h_{w,x,y} \in \mathbb{Z}_{\geq 0}[v]$ by KL positivity.
The trace is therefore:
\[
  \operatorname{tr}_\mu(C'_w)
  \;=\; \sum_{x \in \mathrm{cell}(\mu)} h_{w,x,x}(v)
  \;\in\; \mathbb{Z}_{\geq 0}[v].
\]

\emph{Step~3.}
Combining Steps~1 and~2:
\[
  \operatorname{tr}_\mu(\Pi_q)
  \;=\; v^m \sum_{w} a_w(v) \cdot \operatorname{tr}_\mu(C'_w)
  \;\in\; \mathbb{Z}_{\geq 0}[v].
\]
Since $\operatorname{tr}_\mu(\Pi_q)$ lies in $\mathbb{Z}[q] = \mathbb{Z}[v^2]$
(being the trace of a $\mathbb{Z}[q]$-matrix on a $\mathbb{Z}[q]$-module),
all odd powers of~$v$ vanish.
A polynomial in $\mathbb{Z}_{\geq 0}[v] \cap \mathbb{Z}[v^2]$ has
non-negative coefficients when re-expressed in $q = v^2$.
Hence $\operatorname{tr}_\mu(\Pi_q) \in \mathbb{Z}_{\geq 0}[q]$.
\end{proof}

\section{Eigenvalue positivity}

\begin{theorem}[Eigenvalue Positivity]\label{thm:eigenvalue}
For every $\lambda \vdash n$ and $q > 0$, all nonzero eigenvalues of\/
$\Pi_q$ on $V_\lambda$ are positive.
\end{theorem}

\begin{proof}
\emph{Step~1 (Elementary symmetric functions).}
Let $r_\lambda$ be the rank of $\Pi_q$ on $V_\lambda$ over $\mathbb{Q}(q)$,
and let $e_k(q)$ ($1 \leq k \leq r_\lambda$) be the $k$-th elementary
symmetric function of the nonzero eigenvalues. We have:
\[
  e_k(q) \;=\; \operatorname{tr}\!\Bigl(
    \Pi_q \text{ on } \textstyle\bigwedge^k V_\lambda
  \Bigr)
\]
where $\Pi_q$ acts diagonally on $\bigwedge^k V_\lambda$ via
$\bigwedge^k(\Pi_q|_{V_\lambda})$.
(Here we use the standard identity
$\bigwedge^k(AB) = (\bigwedge^k A)(\bigwedge^k B)$
and the fact that the diagonal $S_n$-action on $\bigwedge^k V_\lambda$
gives $\bigwedge^k \rho(g)$ for each group element~$g$.)

\emph{Step~2 (Decomposition).}
As an $H_q(S_n)$-module,
$\bigwedge^k V_\lambda \cong \bigoplus_\mu V_\mu^{\oplus a_{k,\mu}}$
with $a_{k,\mu} \geq 0$.
By Theorem~\ref{thm:trace}:
\[
  e_k(q) \;=\; \sum_\mu a_{k,\mu} \cdot \operatorname{tr}_\mu(\Pi_q)
  \;\in\; \mathbb{Z}_{\geq 0}[q].
\]

\emph{Step~3 (Non-vanishing).}
Over $\mathbb{Q}(q)$, the matrix $\Pi_q|_{V_\lambda}$ has rank~$r_\lambda$.
Hence $e_k \neq 0$ as a rational function (and therefore as a polynomial,
since $e_k \in \mathbb{Z}[q]$) for $1 \leq k \leq r_\lambda$.
A nonzero polynomial in $\mathbb{Z}_{\geq 0}[q]$ is strictly positive
for all $q > 0$.

\emph{Step~4 (Positivity of all eigenvalues).}
The nonzero eigenvalues $\lambda_1, \ldots, \lambda_{r_\lambda}$ of
$\Pi_q$ on $V_\lambda$ are real for $q > 0$ (as eigenvalues of a real matrix).
They are the roots of the characteristic polynomial of the restriction to the
image:
\[
  p(x) \;=\; \sum_{k=0}^{r_\lambda} (-1)^k\, e_k(q)\, x^{r_\lambda - k},
  \qquad e_0 = 1.
\]
For $x \leq 0$ and $q > 0$: consecutive terms
$(-1)^k\, e_k(q)\, x^{r_\lambda - k}$ and
$(-1)^{k+1}\, e_{k+1}(q)\, x^{r_\lambda - k - 1}$
have the same sign, since the sign change from $(-1)^k \to (-1)^{k+1}$
cancels the sign change from $x^{r_\lambda-k} \to x^{r_\lambda-k-1}$
(as $x \leq 0$).
Hence all terms have a common sign, and $p(x) \neq 0$.
All roots are positive.
\end{proof}

\begin{corollary}[Rank Constancy]\label{cor:rank}
$\operatorname{rank}(\Pi_q|_{V_\lambda})$ is constant for all $q > 0$.
\end{corollary}

\begin{proof}
Since all $r_\lambda$ eigenvalues are positive (hence nonzero) for $q > 0$,
the rank cannot drop below~$r_\lambda$.
And rank cannot exceed~$r_\lambda$ (the generic rank).
\end{proof}

\begin{remark}
The non-negativity argument requires \emph{only} that
$1 + T_s = v \cdot C'_s$ (identifying our factors with KL generators)
and Elias--Williamson positivity.
No case analysis, no explicit eigenvalue computation, no restriction on~$n$.
\end{remark}

\paragraph{Computational verification.}
The trace non-negativity has been verified independently by explicit
computation of eigenvalue polynomials for all $\lambda \vdash n$
with $n \leq 8$.

\begin{thebibliography}{EW14}

\bibitem[EW14]{EW14}
B.~Elias and G.~Williamson,
\emph{The Hodge theory of Soergel bimodules},
Ann.\ of Math.\ \textbf{180} (2014), 1089--1136.

\bibitem[KL79]{KL79}
D.~Kazhdan and G.~Lusztig,
\emph{Representations of Coxeter groups and Hecke algebras},
Invent.\ Math.\ \textbf{53} (1979), 165--184.

\end{thebibliography}

\end{document}
